\begin{table}[htbp]
\centering
\caption{Annualized growth in monthly labor income by educational attainment}
\label{tab:latam_country_annualized_growth}
\small
\begin{tabular}{lrrrrr}
\toprule
Economy & Period & Total & Low & Middle & High \\
\midrule
Argentina & 2016--2023 & -4.5 & -4.6 & -5.2 & -4.8 \\
Bolivia & 2016--2023 & -0.6 & -0.8 & -1.7 & -0.4 \\
Brazil & 2016--2023 & 0.8 & 0.3 & -0.1 & -1.3 \\
Chile & 2016--2023 & 1.4 & 0.9 & -0.8 & -1.8 \\
Colombia & 2016--2023 & 2.7 & -0.6 & -0.6 & -1.2 \\
Costa Rica & 2016--2024 & 0.1 & -0.2 & -1.1 & -1.2 \\
Dominican Republic & 2016--2023 & 0.6 & 0.2 & 1.4 & -0.9 \\
Ecuador & 2016--2024 & -2.3 & -1.8 & -1.3 & -2.9 \\
El Salvador & 2016--2023 & 3.0 & 3.4 & 1.4 & 1.1 \\
Guatemala & 2016--2022 & 0.5 & 2.8 & -2.8 & 0.3 \\
Mexico & 2016--2023 & 1.2 & 1.5 & 0.3 & -0.6 \\
Paraguay & 2017--2023 & -0.1 & 0.8 & -1.1 & -3.7 \\
Peru (Lima and Callao) & 2016--2023 & -1.6 & -1.7 & -1.2 & -2.8 \\
Uruguay & 2016--2023 & 0.0 & -1.3 & -0.4 & 0.2 \\
\bottomrule
\end{tabular}
\begin{minipage}{0.98\textwidth}
\footnotesize
\textit{Note:} Compound annual growth rates, in percent per year, between the initial and final fourth-quarter observations. Total refers to synthetic mean monthly labor income; Low, Middle, and High refer to the mean monthly income of each education group. All underlying income values are measured in 2017 PPP US dollars per employed person per month. Rates use only the two endpoints.
\end{minipage}
\end{table}
